\begin{center}
\refstepcounter{table}\label{tab:local_conformal}
\begin{minipage}{0.92\textwidth}
\textbf{Table \thetable:} Summary of the toxicity-phenotype/MIE-proxy conformal calibration experiment. The point predictor was a deliberately simple structural QSAR model trained on Morgan fingerprints, MW, and logP; docking scores and ADMETLab toxicity-phenotype outputs were used only for local residual calibration. Wilson intervals use the held-out molecules as Bernoulli trials; cluster-bootstrap intervals resample whole Butina test clusters.
\par\vspace{0.5em}
\centering
\begin{tabular}{lr}
\toprule
Quantity & Value \\
\midrule
Model-training molecules & 8,038 \\
Calibration molecules & 2,079 \\
Held-out test molecules & 2,533 \\
\midrule
Global 90\% split-conformal interval width & 1.573 \\
Global 90\% empirical coverage & 89.89\% (2,277/2,533) \\
Global Wilson 95\% CI & 88.66\%--91.01\% \\
Global cluster-bootstrap 95\% CI & 88.29\%--91.37\% \\
Localized 90\% empirical coverage & 89.18\% (2,259/2,533) \\
Localized Wilson 95\% CI & 87.91\%--90.33\% \\
Localized cluster-bootstrap 95\% CI & 87.52\%--90.74\% \\
Localized fallback use & 14/2,533 (0.55\%) \\
Localized calibration coverage (no fallback) & 89.32\% (2,250/2,519) \\
Fallback-subset coverage (diagnostic only) & 64.29\% (9/14) \\
Fallback-subset Wilson 95\% CI & 38.76\%--83.66\% \\
Fallback vs. localized mean width & 1.573 vs. 1.514 \\
Mean local 90\% interval width & 1.515 \\
Median local 90\% interval width & 1.519 \\
Local interval width, 10th--90th percentile & 1.209--1.813 \\
Displayed map range of interval width & 1.20--1.88 \\
\bottomrule
\end{tabular}
\end{minipage}
\end{center}
